%*******************************************************
% Phụ lục: câu hỏi hội đồng có thể đặt ra
%*******************************************************
\chapter{Câu hỏi hội đồng có thể đặt ra}\label{pl:cauhoi}

Phụ lục này dùng để tự luyện trước buổi bảo vệ đề cương. Mỗi câu trả lời chỉ dựa vào những gì đã viết trong chương~\ref{ch:4}, và ghi mục để đọc lại.

\begin{description}[style=nextline, leftmargin=0pt]
\item[Đề tài có giải quyết bài toán số 19 không?] Không trọn vẹn. Đề tài cung cấp thước đo nền cho câu hỏi số 1 và số 3 của bài toán; câu hỏi số 2 và số 4 nằm ngoài phạm vi. Sản phẩm gần nhất với danh mục sản phẩm của bài toán là Chỉ số chất lượng thị trường vốn, có thể là một cấu phần của \enquote{Hệ thống dữ liệu tài chính đô thị} (mục~\ref{sec:lydo}, \ref{sec:phamvi}).

\item[Trung tâm tài chính quốc tế đặt ở Thành phố Hồ Chí Minh và Đà Nẵng, sao lại nghiên cứu cho Hà Nội?] Chính vì vậy. Theo Nghị quyết số 222/2025/QH15, trung tâm tài chính quốc tế đặt tại Thành phố Hồ Chí Minh và Đà Nẵng; Hà Nội không có trong GFCI 40 (Wardle \& Mainelli, 2026) và sắp không còn sàn cổ phiếu niêm yết. Câu hỏi thực tế là Hà Nội nên chuyên môn hóa vào đâu, và câu trả lời cần số đo (mục~\ref{sec:lydo}).

\item[Vì sao không dùng sai số chuẩn phân cụm theo cổ phiếu?] Vì mọi cổ phiếu chuyển sàn cùng một ngày, nên về bản chất chỉ có một lần chuyển sàn; coi 299 mã là 299 quan sát độc lập làm sai số chuẩn quá nhỏ (Conley \& Taber, 2011). Suy luận chính là bootstrap tự hồi quy trên chuỗi chênh lệch theo tháng (mục~\ref{sec:kiemdinh}).

\item[Nếu không bác bỏ được giả thuyết gốc thì sao?] Đề tài báo cáo kích thước hiệu ứng và khoảng tin cậy. Nếu khoảng tin cậy chứa cả 0 lẫn tác động lớn, kết quả là chưa đủ cơ sở kết luận, không phải bằng chứng không có tác động. Hai chỉ số và tệp đăng ký vẫn là sản phẩm dùng được (mục~\ref{sec:tincay}, \ref{sec:lydo}).

\item[Đề tài phát hiện được tác động nhỏ đến đâu?] Với giả định minh họa $\sigma = 0{,}2$ và $\phi = 0{,}6$, logarit Amihud phải giảm khoảng 0,25 trở lên, tức độ phi thanh khoản giảm khoảng 22\%. Giá trị thật sẽ được ước lượng từ giai đoạn tham chiếu và ghi vào tệp đăng ký (mục~\ref{sec:tincay}).

\item[Biên độ và bước giá thay đổi cùng lúc với sàn, làm sao tách ra?] Không tách được hoàn toàn; $\delta$ đo tác động gộp của sàn. Đề tài hạn chế phần thay đổi cơ học bằng ngưỡng lợi suất chung khi tính Amihud và báo cáo theo nhóm mức giá, và ước lượng tác động khác nhau theo mức chịu cú sốc, nhưng các tương tác không đo riêng tác động của từng quy tắc (mục~\ref{sec:giathuyet}, \ref{sec:kiemdinh}).

\item[Sao không dùng tỷ số Sharpe cho khả năng đầu tư?] Vì khi lợi suất vượt trội trung bình âm, danh mục rủi ro hơn lại có tỷ số Sharpe ít âm hơn. Lợi suất tương đương chắc chắn luôn phạt rủi ro cao hơn. Tỷ số Sharpe vẫn được báo cáo là thước đo phụ (mục~\ref{sec:khanangdautu}).

\item[Vì sao cần bảy mô hình, không dùng một?] Để kết luận không phụ thuộc vào sai số của một mô hình. Độ phân tán giữa các mô hình cho biết mức phụ thuộc đó; khoảng cách của \MM{1}, mô hình không cần ước lượng tham số, được báo cáo riêng như phép so sánh sạch nhất (mục~\ref{sec:mohinh}, \ref{sec:khanangdautu}).

\item[\MM{5} là mô hình của chính thí sinh, có thiên lệch không?] Có thể có, nên kết quả được báo cáo cả khi có lẫn khi không có \MM{5}, và mọi mô hình được tinh chỉnh với cùng quy tắc và cùng ngân sách (mục~\ref{sec:mohinh}).

\item[Thí sinh đã dùng dữ liệu HOSE trong các bài trước, làm sao tránh chọn kết quả?] Bằng đăng ký trước trên OSF Registries trước ngày 28/12/2026, khai rõ phạm vi sàn, giai đoạn và chỉ báo đã tính trong các bài trước; mọi thay đổi sau đó được ghi là lệch khỏi kế hoạch (mục~\ref{sec:mohinh}).

\item[FTSE Russell nâng hạng cùng thời gian, có làm hỏng thiết kế không?] FTSE Russell coi cả HOSE và HNX là sàn đủ điều kiện, nên đổi sàn không đổi tư cách sàng lọc. Danh sách đủ điều kiện chỉ gồm mã HOSE và được dùng để loại mã khỏi nhóm đối chứng (mục~\ref{sec:kiemdinh}, \ref{sec:did}).

\item[UPCoM thì sao?] Đợt chuyển ngày 28/12/2026 không gồm UPCoM, nhưng hạn chuyển UPCoM theo thông tư vẫn là hết ngày 31/12/2026. Nếu UPCoM chuyển, nhóm đối chứng UPCoM chỉ được dùng đến tháng liền trước ngày chuyển, và ước lượng chính dựa vào nhóm đối chứng HOSE (mục~\ref{sec:tincay}).

\item[\HH{1} có phải quan hệ nhân quả không?] Không. Điểm GFCI đã chứa dữ liệu giao dịch cổ phiếu, và trung tâm mạnh có thể kéo theo giao dịch cổ phiếu, nên $\beta$ chỉ là liên hệ có điều kiện (mục~\ref{sec:kiemdinh}).

\item[Kết quả giúp gì cho HNX?] Chỉ số chất lượng thị trường vốn có thể được mở rộng sang trái phiếu chính phủ và phái sinh, và báo cáo định hướng chuyên môn hóa hướng tới chính các thị trường HNX tiếp tục vận hành. Thí sinh sẽ xin văn bản xác nhận hoặc cam kết sử dụng kết quả, ưu tiên HNX (mục~\ref{sec:phamvi}).
\end{description}
